\chapter{Conclusiones}

\begin{itemize}
   \item El método empleado en este trabajo permitió estimar los parámetros físicos $i$, $q$, $\Omega_{1}$, $\Omega_{2}$ y $L$ de sistemas binarios eclipsantes a partir de datos fotométricos. Para las binarias ASAS 203113+0513.2 y ASAS 011416+0426.4, los parámetros geométricos $i$ y $q$ resultaron congruentes, dentro de sus incertidumbres, con los valores reportados en estudios previos. Sin embargo, el resto de los parámetros presentó discrepancias superiores al $10\%$ e incluso cercanas al $100\%$ en algunos casos, reflejando algunas limitaciones del método. En el caso de la binaria U Ophiuchi, la ausencia de errores en los datos fotométricos impidió calcular las incertidumbres de sus parámetros; aun así, los valores obtenidos en la etapa de optimización inicial se mantuvieron en concordancia general con los de los demás sistemas analizados.

\item Se logró modelar exitosamente las curvas de luz observacionales utilizando datos fotométricos provenientes de catálogos astronómicos, y el modelo simplificado de Roche permitió reproducir las curvas de luz sintéticas correspondientes. Los valores reducidos del indicador de ajuste SCRM evidencian una buena congruencia entre las curvas observacionales y sintéticas, confirmando que el procedimiento empleado describe correctamente el comportamiento fotométrico de los sistemas estudiados.

\item El algoritmo genético empleado cumplió su función de identificar regiones del espacio de parámetros capaces de reproducir adecuadamente las curvas de luz observacionales. Sin embargo, su carácter inherentemente aleatorio dio lugar, en algunos casos, a soluciones numéricamente diferentes en ejecuciones distintas, lo que indica que, si bien útil, no constituye el método más estable para este tipo de problemas altamente degenerados. Se concluye que los algoritmos genéticos son apropiados para una búsqueda inicial amplia, pero pueden complementarse con métodos deterministas como Simplex para alcanzar soluciones más estables y robustas.

\item El método de Monte Carlo permitió cuantificar la incertidumbre asociada a cada parámetro, incorporando la dispersión de los datos fotométricos. Las distribuciones obtenidas muestran correlaciones y regiones degeneradas del espacio de soluciones. Aunque un mayor número de iteraciones habría permitido obtener estimaciones estadísticamente más precisas y mejor caracterizadas, las limitaciones de tiempo y capacidad computacional impidieron ampliar significativamente el muestreo.

\item El sistema ASAS 061107+1619.5 fue caracterizado por primera vez en este informe. Los parámetros obtenidos presentan incertidumbres moderadas, una distribución relativamente compacta en las simulaciones de Monte Carlo y un valor reducido de SCRM  que confirma que la curva sintética reproduce satisfactoriamente la curva observacional, por lo que esta solución constituye una primera estimación fiable de la geometría del sistema. Aunque no se trata de una solución definitiva, representa un aporte significativo al conocimiento de esta fuente y sienta las bases para futuros estudios observacionales y modelados más avanzados.
\end{itemize}

